\section{Simultaneous finite-place specialization and the denominator budget}
\label{sec:finite}\label{c4:finite-repairs}

The regular twist in Section~\ref{c3:twists} must work for every power from
\(1\) to \(m\) at a given identification \(\iota\).
We prove the specialization result with the order of the roots bounded
\emph{before} the auxiliary primes and their polynomials are chosen.
We also evaluate the sum over primes of the explicit crystalline-loss
allowance. Its application to crystalline classes is a separate arithmetic
hypothesis, specified below.

\subsection{A common point for all the polynomials}
\label{c4:root-spec}
Fix a prime \(\ell\). Let \(v\) be the ordinary additive valuation on
\(\overline{\mathbb Q}_\ell\), with \(v(\ell)=1\), \(v(0)=+\infty\);
write \(\mathcal O\) for its valuation ring. Fix a bound \(s_0\) on the
number of variables, a positive rational \(u\), a nonnegative real
\(\beta\), and \(\epsilon>0\). There exists an integer \(h_0\) such
that the following holds. Let \(P_i\) be any finite collection of
polynomials over \(\mathcal O\) in \(s\le s_0\) variables and let
\(x_i\in\mathcal O^s\) satisfy
\[
 v(x_{ij})\ge u,\qquad \sum_i v(P_i(x_i))\le\beta.
\]
There exist \(\ell^{h_0}\)-th roots of unity \(\zeta_j\), common to all
\(i\), such that \(v(\zeta_j-1)\le u\) and
\[
 v\bigl(P_i((\zeta_j-1)_j)\bigr)\le v(P_i(x_i))+\epsilon
 \quad\text{for every }i .
\]
There is no degree restriction on \(P_i\).

For nonzero \(g=\sum a_k Y^k\in\mathcal O[Y]\), put
\(v_w(g)=\min_k(v(a_k)+kw)\). This is additive on products for
\(w\ge0\): in the product take the sum of the smallest indices
attaining the minima in each factor. Only their product attains the
minimum at that index. Factoring \(g\) into linear factors over
\(\overline{\mathbb Q}_\ell\) now gives
\[
 v_w(g)=v_0(g)+\sum_{v(\alpha)\ge0}\min(w,v(\alpha)),
\]
with roots counted with multiplicity. Indeed each nonintegral root
already contributes \(v(\alpha)\) at \(w=0\) and at \(w\ge0\).
For \(v(y)=w\), all factors except those with \(v(\alpha)=w\)
have \(v(y-\alpha)=\min(w,v(\alpha))\).

Write \(\rho_h=1/\varphi(\ell^h)\), \(h\ge1\). Every primitive
\(\ell^h\)-th root has \(v(\zeta-1)=\rho_h\): its conjugates have
the same valuation and their product, up to sign, is
\(\Phi_{\ell^h}(1)=\ell\). Fix \(0<\delta<1\), \(B\ge0\).
Choose \(h_0\) so large that
\[
 B<\delta\sum_{h\in\mathcal H}\frac1{h+1},\qquad
 \mathcal H=\{1\le h\le h_0:\rho_h\le u\}.
\]
This is possible by divergence of the harmonic series. If
\(v_u(g)\le B\), and \(R_h\) is the multiset of roots with valuation
\(\rho_h\), the distinctness of the \(\rho_h\) and the factor formula
give \(\sum_{h\in\mathcal H}|R_h|\rho_h\le v_u(g)\).
Some \(h\) therefore satisfies \(|R_h|\rho_h(h+1)\le\delta\).

For \(\alpha\in R_h\), apply the same factor formula to
\(F(Y)=(Y+1+\alpha)^{\ell^h}-1\). Its roots include
\(\zeta-1-\alpha\) for every primitive \(\ell^h\)-th root \(\zeta\).
The coefficient of \(Y\) has valuation \(h\) because \(1+\alpha\) is
a unit. Hence
\[
 \sum_{\zeta\ \mathrm{primitive}}\min\{1,v(\zeta-1-\alpha)\}
 \le v_1(F)\le h+1.
\]
Averaging over the \(1/\rho_h\) primitive roots yields one \(\zeta\)
for which the sum over \(\alpha\in R_h\) is at most
\(|R_h|\rho_h(h+1)\le\delta<1\). Every truncated valuation in this
last sum equals the untruncated valuation. Returning to the factors
of \(g\), we obtain
\[
 v(g(\zeta-1))\le v_{\rho_h}(g)+\delta.
\]

Take \(\delta=\min(1,\epsilon)/(s_0+2)\) and use \(B=\beta+1\)
in choosing \(h_0\). Move one coordinate at a time. At a coordinate
still to be moved let \(g_i(Y)\) denote \(P_i\) evaluated at its
current coordinates except for this one. Its old coordinate has
valuation at least \(u\). Apply the one-variable result to
\(g=\prod_i g_i\), whose \(v_u\) is at most the sum of the old
valuations. The excesses
\(v(g_i(\zeta-1))-v_{\rho_h}(g_i)\) are nonnegative and sum to at
most \(\delta\). Since each \(v_{\rho_h}(g_i)\) is at most the
old evaluation, each valuation and their sum increase by at most
\(\delta\) at this step. In particular the sum stays below
\(\beta+1\) at all steps. After at most \(s_0\) steps the total
increase for each evaluation is at most \(\epsilon\). The chosen
\(h_0\) depends only on the prescribed initial bounds.

\paragraph{Powers of one character.}
Suppose instead that \(v(P_i(x_i))\le\beta_i<\infty\), and fix
\(m\ge1\). Apply the preceding argument to the polynomials
\[
 P_{i,j}(Y)=P_i\bigl(((1+Y_a)^j-1)_a\bigr),\qquad 1\le j\le m .
\]
For each original coordinate \(x\), the monic polynomial
\((1+Y)^j-(1+x)\) has \(j\) integral roots whose product is
\(\pm x\). At least one has valuation \(\ge u/j\ge u/m\).
Use this preimage for the starting point of \(P_{i,j}\), the radius
\(u/m\), and the bound \(m\sum_i\beta_i\).
The same tuple of roots now satisfies, for every \(1\le j\le m\),
\[
 v(\zeta_a-1)\le u/m,\qquad
 v\bigl(P_i((\zeta_a^j-1)_a)\bigr)\le v(P_i(x_i))+\epsilon .
\]
If \(u<1\), a root of order \(\ell^h\) so chosen cannot have order
dividing \(j\), since otherwise
\(v(\zeta_a-1)=1/\varphi(\ell^h)\ge1/j>u/m\).
Consequently
\(v(\zeta_a^j-1)=\ell^{v(j)}v(\zeta_a-1)\le j\,u/m\le u\),
where \(v(j)=v_\ell(j)\). This gives nontrivial inertia for every
one of the required powers.

\subsection{The arithmetic assertion needed before specialization}
\label{c4:twist-polynomials}
Here is the exact point at which the Selmer construction enters.
Start with a CM field \(E\), a finite set \(S_0\) containing the places
above \(\ell\), an inertial element \(g_0\) above \(S_0\) away from
\(\ell\), and finite collections of Selmer data and Stickelberger
systems as in Section~\ref{c3:twists}. Normalize module length by
\(\operatorname{length}(\mathcal O'/\ell)=1\) for the coefficient
DVR \(\mathcal O'\). The numerical allowances of the data and systems
are denoted by \(\beta_i\). Assume the coefficient DVR and a torsion level \(M_0\) have been chosen together with
\(M_0>\beta_i+1\) for every Selmer datum and residue field size
exceeding \(2\operatorname{rank}(T_i)
\operatorname{length}_{\rm unscaled}(\mathcal O'/\ell^{M_0})+1\).
The following polynomial existence statement is required.
There is \(s_0\ge1\) such that for every rational \(u>0\) with
\(s_0u<1\), and all sufficiently large \(N\), there are
\begin{itemize}
\item an auxiliary set \(Q\) of at most \(s_0/2\) degree-one primes, with
\(Q,\bar Q,S_0\) disjoint in the sense that \(Q\cup\bar Q\) avoids
\(S_0\) and \(Q\cap\bar Q=\varnothing\);
\item characters \(\lambda_a\) of \(\operatorname{Gal}(\bar E/E)\)
with values in \(\mathbb Z/\ell^N\), induced by
\(z\mapsto z/\bar z\) from tame characters of level \(a\), trivial
above \(\ell\), and surjective on inertia at \(a\);
\item integral polynomials \(P_i\) and points \(x_i\) on these
variables with coordinates of valuation at least \(u\) and
\(v(P_i(x_i))\le\beta_i+s_0u\).
\end{itemize}
For \(\zeta_a^{\ell^N}=1\), put
\(\psi=\prod_{a\in Q}\zeta_a^{\lambda_a}\).
For each system, \(P_i((\zeta_a-1)_a)\) is its twisted
group-ring evaluation at \(\psi\). For each datum with this evaluation
of valuation \(<\beta_i+1\), that valuation bounds the normalized
length of its divisible-coefficient Selmer group twisted by \(\psi\).
Sections~\ref{sec:selmer-polynomial-descent} and~\ref{sec:terminal-selmer} give a polynomial construction for abstract Selmer and Stickelberger data, with precise hypotheses on the local conditions, scalar inertia and Euler factors. Application to the critical pairs requires those hypotheses for their actual representations and values.

To deduce the general twisting assertion, it suffices to take
\(0<\epsilon<1\), decreasing the prescribed tolerance if necessary.
Choose \(s_0u\le\epsilon/2\), then \(h_0\) for precision
\(\epsilon/2\), powers \(1,\ldots,m\) and bounds \(\beta_i+1\).
Only now choose \(N\ge h_0\) and the corresponding \(Q\) and
polynomials. Every specialized evaluation at \(\psi^j\) has valuation
\(\le\beta_i+\epsilon<\beta_i+1\). At \(a\), all characters but
\(\lambda_a\) are unramified and a generator of tame inertia has
\(\lambda_a\)-value prime to \(\ell\), so its image under \(\psi^j\)
has the order of \(\zeta_a^j\). The conjugate prime has the same
order by anticyclotomicity. Thus the sum of the valuations of these
inertial values minus one is at most \(2|Q|u\le\epsilon/2\).
This proves the required simultaneous specialization, conditional on
the polynomial existence statement. The application to lifts further
requires the asserted Galois representations, local conditions and
identification of the critical ratios with the Stickelberger systems.

\subsection{The central average for the rank-one line}
\label{c4:central-average}
Fix one rank-one lift datum and let \(\phi\) be the lift of the
constant source form. Put \(P=P_{\mathbb C\phi}\); the case
\(\phi=0\) is immediate. The test-form condition says that a form is
integral at every place above \(\ell\) and orthogonal to the false
lifts lying in the kernel of the normalized descent. The clearing
operator on \(\mathbb C\phi\) applies to forms orthogonal to false
lifts lying in \((\mathbb C\phi)^\perp\). These conditions are
connected by a class sum.

Let \(s_c\) run over representatives of the \(N_0\) chosen-level
classes on \(U(1)\), integral above \(\ell\). To arrange integrality,
write a representative as \(a/\bar a\) as below and use weak
approximation in \(E\) to make \(a\) a unit at all places above
\(\ell\) by multiplication by a rational idele from \(E^\times\).
Let \(T(z)\) be central translation on the target multiplied by
the unitary descent character at \(z\); \(T(z)^*=T(z^{-1})\).
For the lift kernel \(K(g)\) the central action is
\(T(z^{-1})K(g)=K(gz)\). This uses the diagonal scalar action in the
Weil representation (see~\cite{weil} for the representation): its character on the defining block is the
inverse of the factor multiplying the translation. We use here the
specified Weil model and its scalar action. Since the stabilizers
are trivial, \(\phi=\sum_c K(s_c)\), and translating all representatives
permutes the classes. Thus \(T(z)\phi=\phi\) and, for every \(g\),
\(\sum_c T(s_c^{-1})K(g)=\phi\).

Write \(S=\sum_c T(s_c)\). If \(z\perp\phi\), the descent of
\(Sz\) vanishes: its value at \(g\) is
\(\langle K(g),Sz\rangle=\langle\phi,z\rangle=0\), with the
usual kernel realization of the normalized descent. Central
translations commute with each finite Hecke operator and with the
residue-type idempotent because both are sums of translations.
Thus \(S\) carries a false lift in \((\mathbb C\phi)^\perp\)
at a fixed place into a false lift in the descent kernel at the same
place. If \(\Xi'\) is a test form, \(S^*\Xi'\) is therefore
orthogonal to the former false lifts. It is also integral at
\(\iota\): the representatives and their inverses are integral
above \(\ell\), so scalar translation preserves the local Weyl
lattice, and the normalizing character is an \(\ell\)-adic
unit there (the finite-order part is a root of unity and the
algebraic-weight factors are local units).
Finally \(PS^*\Xi'=N_0P\Xi'\), since \(S\phi=N_0\phi\).

Suppose the rank-one clearing estimate with exponent \(\gamma\)
has been established. Its operator on the line is a scalar \(\mu\)
with \(|\iota\mu|\ge e^{-\gamma}\); its pairing clause gives
\(|\iota\langle \Xi,\mu P S^*\Xi'\rangle|\le1\) for any form
\(\Xi\) integral at all places above \(\ell\).
Since \(\langle P\Xi,P\Xi'\rangle=\langle\Xi,P\Xi'\rangle\),
\[
 |\iota\langle P\Xi,P\Xi'\rangle|
 \le e^\gamma\,|\iota N_0|^{-1}
 =\exp\bigl(\gamma+v_\ell(N_0)\log\ell\bigr).
\]
Only one clearing scalar occurs. Applying the two projections
separately would lose this coefficient.

\subsection{Summing the explicit loss}
\label{c4:defect-sum}
Write \(h=r^8\), \(M=p_0p_1\), and let \(\Pi\) be the product of the
positive integers in \((b-30h^3,b+30h^3]\), with \(b\) identified
with its integer value. The local allowance for crystalline defects is
\[
 \lambda_\ell=4\cdot10^6 f(m/k)
     \left[bv_\ell(M)+v_\ell\bigl((2b)!\,\Pi^{h^3}\bigr)\right].
\]
For every positive integer \(A\) and finite prime set \(F\) unique
factorization gives \(\sum_{\ell\in F}v_\ell(A)\log\ell\le\log A\).
The interval has at most \(60h^3\) members. Therefore
\begin{align*}
 \sum_{\ell\in F}\lambda_\ell\log\ell
 &\le 4\cdot10^6 f(m/k)
   \bigl[b\log M+2b\log(2b)+60h^6\log(b+30h^3)\bigr]\\
 &\le 4\cdot10^6 e^{180} f(m/k)bL .
\end{align*}
Indeed \(p_i\le b^{11}\), \(r\le3x^5\le x^6\), and
\(b\ge x^{10^{12}}\), so \(60h^6\le b\), \(30h^3\le b\) and
\(\log(2b)\le2L\). The bracket is already at most \(28bL\).
This computes the sum of the proposed local allowances. The
comparison of crystalline cocycles with these allowances for the
main summands enters the generator-wise congruence estimate.

\paragraph{Field and class counts.}
The elementary field-size bounds can be taken directly from the
parameterization. With tower exponent \(s\ge1\), the root order is
\(q_1=q_0\,3^{\lfloor\log_3s\rfloor+5}2^s\), where \(2^s=m^{22}\),
\(s\le32\log x\), \(m\le2\sqrt{x}\), and \(q_0,d\le\log x\).
Thus \(q_1\le 2^{36}x^{11}(\log x)^2\le x^{23/2}\); the latter
inequality follows from \(\log x\ge100\log L\) and
\(\log L\ge64\cdot10^{12}\).
The field \(E=\mathbb Q(\sqrt{-d},\zeta_{q_1})\) has
degree \(n_E=2f\le2q_1\) and embeds in the cyclotomic field of
conductor \(C=q_1d\le x^{13}\). Here the quadratic field is contained
in the \(d\)-th cyclotomic field by the quadratic Gauss sum
(\(d\equiv3\bmod4\)). The discriminant formula for cyclotomic fields
gives root discriminant at most \(C\); the relative discriminant
formula gives the same upper bound for the subfield \(E\).

Minkowski's ideal-class bound gives an integral ideal of norm
\(k_0\le B=2^{n_E}|\operatorname{disc}E|\) in every class.
An ideal \(I\) of norm \(k_0\) contains \((k_0)\); the ideal
\(I/(k_0)\) in \(\mathcal O_E/(k_0)\) is principal, since this
ring is a product of quotients of DVRs. Hence \(I=(k_0,a)\) and
there are at most \(k_0^{n_E}\) choices of \(a\bmod k_0\). We obtain
\[
 \log h_E\le(n_E+1)n_E\log(2C)\le x^{28}.
\]
Let \(N_0\) be the number of chosen-level classes of the rank-one
unitary group of \(E/E^+\). The map from finite ideles of \(E\)
to the norm-one finite ideles given by \(z\mapsto z/\bar z\) is
onto. Locally this is Hilbert's Theorem~90~\cite[II, Corollary~1.23]{milneCFT} (at a split pair, take
\((z,1)\)). At an unramified nonsplit place one can make an antecedent
a unit by multiplying by a power of a base-field uniformizer,
so one obtains a finite idele. For the quadratic local assertion
one may explicitly take the antecedent \(1+y\) for a norm-one
\(y\ne-1\), and a nonzero element of trace zero for \(y=-1\).
The image of \(E^\times\) is rational and an integral unit idele
congruent to \(1\) modulo \(M\mathcal O_E\) maps into the chosen
level: it reduces to \(1\) at each chosen place above \(p_0,p_1\).
Thus the ray class group modulo \(M\mathcal O_E\) surjects onto
the indicated unitary class quotient. Its order is at most
\(h_E|(\mathcal O_E/M)^\times|\le h_E M^{n_E}\). In particular
\[
 \log N_0\le x^{28}+22 n_E L\le x^{30}.
\]
Finally, at the chosen places the group defining \(G_0\) is a
product of \(\operatorname{GL}_{rt}\) over residue fields,
\(t=5r\). At all chosen places over \(p_i\) the sum of residue
degrees is at most \(n_E\). Counting all matrices gives
\(\log G_0\le (rt)^2 n_E\log M\le x^{32}\), since
\(rt\le45x^{10}\), \(n_E\le2x^{23/2}\), and \(L\le x^{1/100}\).
These estimates also apply to every finite sum of prime valuations
of the three integers with logarithmic weights.

\paragraph{A uniform cohomological integer.}
The first-order estimate uses one positive integer \(O_*\) whose
normalized prime valuations control odd Tate-twist classes at all
torsion levels for \(0\le e<t\), for coefficient DVRs containing the roots of unity of order \(\exp((\mathbb Z/C\mathbb Z)^\times)\), with the unramified classes when
\(e=0\). The required bound is
\(O_*\le 2^{4(C+t+3)^{27}}h_E\); Section~\ref{sec:odd-cocycles-reduction} gives a classical derivation for the specified cocycle groups. With this bound, \(\log O_*\le x^{2000}\)
follows from \(C\le x^{13}\), \(t\le15x^5\), and the preceding
class-number estimate. The denominator-error valuation is exactly
\[
 t v_\ell(O_*)+
 \{8ft+2+(4t+1)h^{22}+h\}v_\ell(G_0).
\]
As \(h\le x^{48}\), the coefficient in braces is at most \(x^{1100}\);
the sum of this error with logarithmic prime weights is therefore
at most \(x^{4000}\). The twist error
\(u_\ell=(4f+3)r v_\ell(G_0)+2^{-\ell}\) has weighted sum
at most \(x^{999}\) by the group-order estimate and
\(\sum_{n\ge1} n2^{-n}=2\).

Adding the defect allowance, \((1+2t)u_\ell\), the denominator
error, \(v_\ell(G_0)\) and \(v_\ell(N_0)\) gives a sum with
logarithmic weights at most \(B_0=fxbL^3/(mk)\).
Indeed the defect term divided by \(B_0\) is at most
\(16\cdot10^6e^{180}/L^2\le 1/2\), by \(m^2\le4x\).
All other terms sum to at most \(x^{4002}\le B_0/2\), since
\(1\le f\), \(k\le r+1\le4x^5\), \(m\le2\sqrt{x}\), and
\(b\ge x^{10^{12}}\). This closes the numerical estimate for
the local denominators with the asserted arithmetic inputs.

\subsection{A uniform integer for odd cyclotomic cocycles}
\label{sec:odd-cocycles-reduction}
Write $C=q_1d$, $Y=C+t+3$ and $G_F=\operatorname{Gal}(\overline{\mathbb Q}/F)$.
Let $\chi_\ell$ be the cyclotomic character and $c$ a fixed complex conjugation.
For a finite extension of $\mathbb Q_\ell$ with integer ring $\mathcal O$ put
$A_n=\mathcal O/\ell^n\mathcal O$.
We use the integral valuation $v_{\mathcal O}$ with $v_{\mathcal O}(\pi)=1$.
The module $Z_e(E,n)$ consists of continuous cocycles for $A_n(-e)$
which vanish on inertia away from $\ell$. For $e>0$ they satisfy
$z(cgc^{-1})=-z(g)$ and vanish on the kernel of $\chi_\ell$ in every
decomposition group above $\ell$. For $e=0$ they vanish on inertia
above $\ell$ as well and there is no parity requirement. These are
cocycles before quotienting by coboundaries. We establish the uniform bound in Appendix D using ordinary
cyclotomic class groups at small primes and the refined Coates--Sinnott
theorem at large primes.

\paragraph{A common integer from Bernoulli polynomials.}
For $1\le e<t$ write $r=e+1$, $H=(\mathbb Z/C\mathbb Z)^\times$,
$D_r=(r+2)!(r+1)$, and put
\[
 b(a)=D_r C^r B_r(a/C),\qquad 0\le a<C.
\]
Here $B_j(X)$ is the Bernoulli polynomial with $B_1(0)=-1/2$.
The polynomial $P(X)=(r+2)!C^{r+1} B_{r+1}(X/C)$ has integral
coefficients. Indeed $(j+1)!B_j(0)$ is integral by induction in the
Bernoulli recursion $(j+1)B_j(0)=-\sum_{i<j}{j+1\choose i}B_i(0)$ for $j\ge1$.
Also $P'(a)=b(a)$ and $P(X+C)-P(X)=D_r C X^r$. Telescoping and using
the integral Taylor expansion modulo $(Cm)^2$ give, for each $m\ge1$,
\[
 D_r\sum_{i=0}^{m-1}(a+iC)^r\equiv m b(a)\pmod{m^2}.
\]
For $u\equiv1\bmod C$ prime to $m$, multiplication by $u$ permutes residues
modulo $Cm$ and preserves every character $\psi$ modulo $C$. Replacing
$ua$ by $ua-Cm\lfloor ua/(Cm)\rfloor$ in this permutation, expanding the
$r$th power modulo $(Cm)^2$, multiplying by $D_r$ and using the last
congruence residue class by residue class gives
\[
 (u^r-1)\sum_{a\bmod C}\psi(a)b(a)
 \equiv D_r r C u^{r-1}
 \sum_{a=0}^{Cm-1}\Big\lfloor\frac{ua}{Cm}\Big\rfloor a^{r-1}\psi(a)
 \pmod m .
\]
The division by $m$ here takes place after a congruence modulo $m^2$
in the ring of character values.

Define $g(a)=2b(a)+(-1)^e{\bf1}_{a=1}+{\bf1}_{a=-1}$ for $a\in H$
and $T=(g(v/w))_{w,v\in H}$, using least nonnegative representatives.
Its eigenvalue on a character $\psi$ is
$\lambda_\psi=2\sum_{a\in H}\psi(a)b(a)+(-1)^e+\psi(-1)$.
The identity $\zeta(1-r,x)=-B_r(x)/r$ expresses the sum as
$-D_r C r L(-e,\psi)$. This value vanishes exactly for
$\psi(-1)=(-1)^e$. Indeed the functional equation for the primitive
inducing character reduces nonvanishing in the other parity to the
Euler product at $1+e>1$, and the extra Euler factors at $-e$ are nonzero.
Thus $\det T\ne0$. The adjugate identity evaluated at the identity
entry of the character vector shows $\lambda_\psi\mid\det T$ in the
ring of character values, including at primes dividing $|H|$.

Set
\[
 X_e=|\det T|\prod_{u\in\{C+1,\,2C+1\}}
              (u^e-1)(u^{e+1}-1).
\]
The two choices of $u$ are coprime; hence at each prime $\ell$ there
is a choice prime to $\ell$ for which both factors jointly divide $X_e$.
If $\psi(-1)=-(-1)^e$, choose $v$ with $X_e=v\lambda_\psi$.
For every $u\equiv1\bmod C$ and every $m\ge1$ coprime to $u$, the resulting congruence is
\[
 (u^{e+1}-1) X_e\equiv
 2w\sum_{a=0}^{Cm-1}\Big\lfloor\frac{ua}{Cm}\Big\rfloor a^e\psi(a)
 \pmod m,\qquad w=vD_r r C u^e .
\]
The multiplier $w$ is independent of $m$.

Here $1\le X_e\le 2^{Y^6}$. In fact the Fourier series of $B_r$ for
$r\ge2$ gives $|B_r(x)|\le 2r!\zeta(r)/(2\pi)^r\le r!$ on $[0,1]$.
As $r+2\le Y$ and $C\le Y$, the entries of $T$ have modulus at most
$4Y^{3Y+1}$. The Leibniz formula bounds its determinant by
$Y^Y(4Y^{3Y+1})^Y$; the four extra factors contribute at most
$(2Y+1)^{4r}$. Their product has binary logarithm below $Y^6$
for $Y\ge 8$, as holds here.

\paragraph{Annihilation by classical Stickelberger reciprocity.}
For each $\psi:H\to\mathcal O^\times$ write $\alpha=\chi_\ell^e\psi$,
inflating $\psi$ to $G_{\mathbb Q}$. Let $V_\psi$ be the continuous
cocycles on $G_{\mathbb Q}$ for $\alpha^{-1}$ with values in $A_n$
which vanish at $c$ and at those elements of $G_{\mathbb Q(\mu_C)}$
lying in inertia away from $\ell$ or in a decomposition group at
$\ell$ with trivial cyclotomic character. Suppose $\alpha(c)=-1$.
We claim $2X_e^2 V_\psi=0$. We may take $n\ge1$.
Set $F=C\ell^n$, $K'=\mathbb Q(\mu_F)$, and choose
$u\in\{1+C,1+2C\}$ coprime to $\ell$.
If $z\in V_\psi$, the restriction $z|_{G_{K'}}$ is a homomorphism:
$\alpha$ is trivial on $G_{K'}$ modulo $\ell^n$. It vanishes on
inertia outside $\ell$. At a place above $\ell$, choose $\tau$ in
local inertia with $\chi_\ell(\tau)=u$. This is possible because the
local $\ell$-power cyclotomic extension is totally ramified with full
group $\mathbb Z_\ell^\times$. Inertia fixes prime-to-$\ell$ roots of
unity and $u\equiv1\bmod C$, hence $\psi(\tau)=1$.
If $h$ belongs to the local decomposition subgroup of $G_{K'}$,
then the commutator $\tau h\tau^{-1}h^{-1}$ has cyclotomic character one
and fixes $\mu_C$. Its $z$-value is zero by the stated local condition;
the cocycle rule evaluates that value as $(u^{-e}-1)z(h)$ in $A_n$.
It follows that $y=(u^e-1)z|_{G_{K'}}$ factors by global reciprocity
through the ordinary class group $\operatorname{Cl}(K')$.
No archimedean modulus occurs because $K'$ is totally imaginary.

Put $\theta_F=\sum_{b\in(\mathbb Z/F)^\times}(b/F)\sigma_b^{-1}$,
where $0<b<F$ and $\sigma_b$ raises roots of unity to the $b$th power.
The classical Stickelberger theorem annihilates $\operatorname{Cl}(K')$
by the integral group-ring element
\[
 (u-\sigma_u)\theta_F=\sum_{b\in(\mathbb Z/F)^\times}
                 \lfloor ub/F\rfloor\,\sigma_b^{-1}
\]
(L.~C.~Washington, \emph{Introduction to Cyclotomic Fields}, 2nd ed.,
Thm.~6.10~\cite{washington}). The extension is abelian over $\mathbb Q$ and $(u,F)=1$,
as required by that theorem. Conjugation by a lift of $\sigma_b^{-1}$
on $G_{K'}$ multiplies the value under $y$ by $b^e\psi(b)$ in $A_n$;
this is the cocycle identity and the fact that the conjugated element
acts trivially. Reciprocity is equivariant for field automorphisms.
Consequently $U_n y=0$, where $U_n$ denotes the sum of
$\lfloor ub/F\rfloor b^e\psi(b)$ over the units modulo $F$.

Write $T_i$ for the full floor sum at modulus $C\ell^i$. Terms divisible
by $\ell$ give $\ell^e\psi(\ell)T_{n-1}$; terms noncoprime to $C$ vanish.
Thus $U_n=T_n-\ell^e\psi(\ell)T_{n-1}$. Apply the Bernoulli congruence
at $m=\ell^n$ and $\ell^{n-1}$ with the same $w$. Since $e\ge1$ it gives
\[
 2wU_n\equiv(1-\ell^e\psi(\ell))(u^{e+1}-1)X_e\pmod{\ell^n}.
\]
The first factor on the right is a unit. Hence
$(u^{e+1}-1)X_e(u^e-1)z|_{G_{K'}}=0$. Since
$(u^{e+1}-1)(u^e-1)\mid X_e$, the cocycle $X_e^2 z$ is zero
on $G_{K'}$. It therefore descends to the abelian group
$\operatorname{Gal}(K'/\mathbb Q)$. Comparing its values on $gc$ and
$cg$, using $z(c)=0$ and $\alpha(c)=-1$, gives $2 X_e^2z=0$.

\paragraph{Integral character descent.}
Suppose $\mathcal O$ contains the roots of unity of order
$\exp H$. Choose lifts $\sigma_a\in G_{\mathbb Q}$ for $a\in H$ with
$\sigma_1=1$, and put $\delta(g)=cgc^{-1}g^{-1}$. All conjugates
of $\delta(g)$ belong to $G_{\mathbb Q(\mu_C)}\subset G_E$ and have
trivial cyclotomic character. For $z\in Z_e(E,n)$ set
\[
 z_\psi(g)=\sum_{a\in H}\alpha(\sigma_a)\,
       z(\sigma_a\delta(g)\sigma_a^{-1}).
\]
The summand does not depend on the choice of lift: replacing $\sigma_a$
by $h\sigma_a$ for $h\in G_{\mathbb Q(\mu_C)}$ multiplies the $z$-value
by $\chi_\ell(h)^{-e}$ and its scalar by $\chi_\ell(h)^e$.
This uses the cocycle identity for conjugation of an element acting
trivially on the coefficient module. Now
$\delta(gg')=\delta(g)\,g\delta(g')g^{-1}$.
Right multiplication permutes the cosets of the lifts, so
$z_\psi(gg')=z_\psi(g)+\alpha(g)^{-1}z_\psi(g')$.
It vanishes at $c$. If $g$ belongs to one of the designated local
sets in $G_{\mathbb Q(\mu_C)}$, both its conjugates by $\sigma_a$
and by $\sigma_a c$ lie in the corresponding local sets of $G_E$,
so their quotient has zero cocycle value. Hence $z_\psi\in V_\psi$.

For $h\in G_E$ orthogonality of the characters gives
\[
 \sum_\psi z_\psi(h)=|H|z(\delta(h))
              =|H|\{z(chc^{-1})-z(h)\}=-2|H|z(h).
\]
If all odd components vanish, the left side is invariant under
$h\mapsto chc^{-1}$ because each even component is zero at $c$ and
has $\alpha(c)=1$. Consequently $4|H|z=0$.
Applying this fact to $2X_e^2 z$ proves in particular that $B_e=16|H|X_e^2$
annihilates $Z_e(E,n)$ for all $\ell$, including $\ell=2$.

\paragraph{A length bound at large primes by global duality.}
Fix $\ell>Y^2$ and $e>0$, and set $j=e+1$, $\varepsilon=(-1)^j$.
First take $\mathcal O=\mathbb Z_\ell$. Let $S_0$ consist of the
$\ell$-adic and archimedean places of $E$.
Write $M_n=\mathbb Z/\ell^n(-e)$ and $M_n^D=\mathbb Z/\ell^n(j)$
for its Cartier dual. Passage to cohomology gives an injection
\[
 Z_e(E,n)\lhook\joinrel\longrightarrow
 \ker\left[H^1(G_{E,S_0},M_n)\longrightarrow
           \bigoplus_{v\mid\ell}H^1(E_v,M_n)\right]^{c=\varepsilon}.
\]
Indeed coboundaries are invariant under precomposition by
$g\mapsto cgc^{-1}$, whereas our cocycles change sign, and $2$
is invertible. The inertia conditions give descent to $G_{E,S_0}$
because the coefficient action is unramified away from $\ell$.
The natural action of $c$ on cohomology includes its action
$(-1)^e$ on coefficients, yielding the displayed parity. At
$v\mid\ell$ a cocycle factors through the image of $\chi_\ell$ by
the local vanishing condition. This image is $\mathbb Z_\ell^\times$
because $E_v$ is unramified over $\mathbb Q_\ell$ ($\ell>C$).
On this abelian group, commutation with an element $b$ for which
$b^{-e}-1$ is a unit makes every cocycle a coboundary: equating
$z(bg)$ and $z(gb)$ gives $(b^{-e}-1)z(g)=(\chi_\ell(g)^{-e}-1)z(b)$.
Such $b$ exists since $0<e<\ell-1$. Complex places contribute zero.

Poitou--Tate duality for finite modules pairs this localization
kernel perfectly with the localization kernel in
$H^2(G_{E,S_0},M_n^D)$; the pairing respects field automorphisms,
so the same sign occurs on the dual. In particular
\[
 \operatorname{length}_{\mathbb Z_\ell}Z_e(E,n)
 \le \operatorname{length}_{\mathbb Z_\ell}
       H^2(G_{E,S_0},\mathbb Z_\ell(j))^{c=\varepsilon}.
\]
Here the integral group surjects onto finite-level $H^2$:
the finite-module groups have $\ell$-cohomological dimension at most
two in this setting, and the finite cohomology groups in the inverse
system satisfy the Mittag--Leffler condition. Applying the sign
idempotent preserves surjectivity.

Let $S$ consist of $\infty$ and the primes of $\mathbb Q$ dividing $C$,
and write $U=\operatorname{Spec}\mathcal O_{E,S}[1/\ell]$.
The integral group just used injects into
$H^2_{\mathrm{\acute et}}(U,\mathbb Z_\ell(j))$. Indeed Galois cohomology
with ramification allowed at $S_0$ computes the cohomology of
$\operatorname{Spec}\mathcal O_E[1/\ell]$ for these coefficients.
By the localization sequence and purity the terms on the left when
the places $v\mid C$ are removed are $H^0(k_v,\mathbb Z_\ell(j-1))$.
They vanish because $|k_v|^{j-1}-1\ne0$.
These uses of finite-module duality and localization are the usual
Poitou--Tate theorem and the localization sequence for an arithmetic
curve (see J.\,S.~Milne, \emph{Arithmetic Duality Theorems}, 2nd ed.,
Parts I and II~\cite{milneADT}).

Apply the refined Coates--Sinnott theorem of C.~Greither and
C.\,D.~Popescu, \emph{An equivariant main conjecture in Iwasawa
theory and applications}, J.~Algebraic Geom., Thm.~6.11~\cite{gp}, to
$E/\mathbb Q,S,\ell,j$. This is an abelian CM extension of the
totally real field $\mathbb Q$, $j\ge2$, $S$ contains every ramified
prime, and $\ell$ is odd. The required $\mu$-invariant vanishes by
the Ferrero--Washington theorem~\cite{fw} since $E(\mu_\ell)$ is abelian over
$\mathbb Q$. Further, $H^1_{\mathrm{\acute et}}(U,\mathbb Z_\ell(j))$
has zero $\ell$-torsion by the coefficient exact sequence: already
$H^0(E,\mathbb Z/\ell(j))=0$ because $0<j<\ell-1$ and local inertia
maps onto $(\mathbb Z/\ell)^\times$. The theorem therefore identifies
the zeroth Fitting ideal of the $\varepsilon$-part of $H^2$ with the
ideal generated on that part by $\Theta_{E/\mathbb Q,S}(1-j)$.
Indeed its sign idempotent for this base field is
$(1+(-1)^jc)/2$.

To read its length, extend scalars to the integers $\mathcal R$ of
an unramified splitting field of $\operatorname{Gal}(E/\mathbb Q)$
over $\mathbb Q_\ell$. Such a field suffices since the group order
divides $\varphi(C)<\ell$. Its group ring is a product of copies of
$\mathcal R$. Base change preserves Fitting ideals, and the normalized
length of the extended module equals its original $\mathbb Z_\ell$-length.
Over each discrete valuation ring the valuation of the Fitting
generator is the length. Thus the required upper bound is the sum of
the valuations of $L_S(-e,\psi^{-1})$ over the characters of
$\operatorname{Gal}(E/\mathbb Q)$ of sign $\varepsilon$. Inflating
them to $H$ gives exactly the $L$-functions at modulus $C$, since
$S$ removes every prime dividing $C$. For the relevant parity their
values equal the eigenvalues $\lambda_{\psi^{-1}}$ of $T$ divided by
$-2D_r C r$, an $\ell$-adic unit. All other eigenvalues of the
integer matrix $T$ are $\ell$-integral. Consequently
\[
 \operatorname{length}_{\mathbb Z_\ell}Z_e(E,n)
 \le v_\ell(\det T)\le v_\ell(X_e).
\]
For general $\mathcal O$ the cocycle equations and local and parity
conditions hold coordinatewise in a $\mathbb Z_\ell$-basis of
$\mathcal O$. Hence $Z_e(E,n;\mathcal O)$ is the scalar extension
of $Z_e(E,n;\mathbb Z_\ell)$, and the length bound becomes
$v_{\mathcal O}(X_e)$, including for ramified coefficient rings.

\paragraph{Lengths at small primes.}
For $\ell\le Y^2$ a uniform detecting set comes from class field theory.
Put $M=C\ell^2$, $K=\mathbb Q(\mu_M)$, $d_K=[K:\mathbb Q]\le M$.
The maximal exponent-$\ell$ abelian quotient of $G_K$ unramified away
from $\ell$ has $\mathbb F_\ell$-rank at most
\[
 \log_2 h_K+\sum_{v\mid\ell}([K_v:\mathbb Q_\ell]+1)
 \ \le\ \log_2 h_K+2M .
\]
Indeed global reciprocity identifies the quotient by the local inertia
images with a quotient of the ideal class group, and local reciprocity
maps the local units onto abelian inertia. Their pro-$\ell$ factors
have $\mathbb Z_\ell$-rank $[K_v:\mathbb Q_\ell]$ and cyclic torsion,
also at $\ell=2$. Choose lifts of generators and adjoin representatives
for $G_E/G_K$, at most $M$ more points.
For $z$ killed by $\ell$, restriction to $G_K$ is a homomorphism
since $K$ contains $\mu_\ell$. Thus vanishing at the points forces
$z=0$. For a general $z$ vanishing there, induct on its exponent,
first applying induction to $\ell z$ and then the exponent-$\ell$ case.
This works for ramified coefficient rings because $\ell z=0$ makes
the twist trivial on every value over $G_K$.

For completeness, Minkowski's bound and $|\operatorname{disc}K|\le M^{d_K}$
give an integral representative of norm $j\le(2M)^{d_K}$ for each
ideal class. Every ideal of norm $j$ contains $j\mathcal O_K$ and
is generated by $j$ and one other element, since the quotient by
$j\mathcal O_K$ is a product of principal local rings. There are at most
$j^{d_K}$ choices. It follows that $\log_2 h_K\le
d_K(d_K+1)\log_2(2M)$. Since $M\le Y^5$, the total number of points
is at most $Y^{20}$ (use $\log_2(2M)\le M$ for $M\ge2$ and $M\ge4$).
The class field theory used here is the idelic reciprocity theorem
and the Hilbert class field theorem; see J.\,S.~Milne,
\emph{Class Field Theory}, chapters V and VIII~\cite{milneCFT}.

At small primes evaluation thus embeds $Z_e(E,n)$ in at most
$Y^{20}$ copies of $A_n[B_e]$. Each copy has length at most
$v_{\mathcal O}(B_e)$. At large primes the preceding stronger
bound $v_{\mathcal O}(X_e)$ is at most $Y^{20}v_{\mathcal O}(B_e)$.
For $e=0$ class field theory identifies
$Z_0(E,n)$ with $\operatorname{Hom}(\operatorname{Cl}(E),A_n)$;
decomposing the class group into cyclic summands bounds its length by
$v_{\mathcal O}(h_E)$. Therefore
\[
 O_*=h_E\left(\prod_{1\le e<t}16\varphi(C)X_e^2\right)^{Y^{20}}
 \quad\text{satisfies}\quad
 \operatorname{length}_{\mathcal O} Z_e(E,n)\le v_{\mathcal O}(O_*)
\]
for every coefficient ring with the stated roots and every torsion level. Since each factor is
at most $2^{4Y^6}$ and there are at most $Y$ factors,
$O_*\le h_E2^{4Y^{27}}$. Dividing by $v_{\mathcal O}(\ell)$
gives the normalized valuation convention used in the denominator
estimate. This calculation makes the integer independent of
$\ell,\mathcal O,e,n$ and retains the factor $h_E$ only once.

